\chapter{反常积分}
\label{chap:improper-integrals}

Riemann 积分只在有限闭区间上处理有界函数。把区间端点移到无穷远，
或允许函数在端点、内点附近无界，都超出了这个定义。反常积分不另造一种
有限和，而是把每个截断后的 Riemann 积分组成一个极限问题。

以下各判别均以被积函数在远离奇点的每个有限闭子区间上 Riemann 可积
为前提。仅有绝对值可积或点态比较，不能代替这一局部可积性假设。
需要连续性或连续可导性的定理另行注明。

\section{无穷区间、无界函数和多奇点积分}

\begin{definition}[无穷区间上的积分]
设 \(f\) 在每个 \([a,R]\) 上 Riemann 可积。若有限极限
\[
 \lim_{R\to+\infty}\int_a^R f(x)\,dx
\]
存在，则称 \(\int_a^{+\infty}f\) \term{收敛}，并以该极限为其值；
否则称为发散。
\end{definition}

对 \((-\infty,b]\) 作相反方向的截断。双无穷区间必须在某个有限点
\(c\) 分成两段：
\[
 \int_{-\infty}^{+\infty}f
 :=
 \int_{-\infty}^c f+\int_c^{+\infty}f,
\]
且两段都收敛时才称原积分收敛。只考察对称极限
\(\lim_{R\to\infty}\int_{-R}^R f\) 得到的是 Cauchy 主值，不能代替
上述定义。例如 \(f(x)=x\) 的对称主值为零，但两侧反常积分都发散。

\begin{definition}[无界点附近的积分]
设 \(f\) 在每个 \([a,b-\eta]\) 上可积。若
\[
 \lim_{\eta\downarrow0}\int_a^{b-\eta}f(x)\,dx
\]
存在且有限，就定义 \(\int_a^b f\) 为该极限。左端奇点类似。若
\(c\in(a,b)\) 是奇点，则必须分别要求
\[
 \int_a^c f,\qquad \int_c^b f
\]
收敛；不能让两侧发散量互相抵消。
\end{definition}

若有有限多个奇点，就在每个奇点处分段并逐段要求收敛。这个规则使积分值
与截断速度无关。

\begin{example}[\(p\)-型积分]
\label{ex:p-improper-integrals}
直接计算截断积分可得
\[
 \lim_{R\to\infty}\int_1^R \frac{dx}{x^p}
 =
 \begin{cases}
 \dfrac{1}{p-1},&p>1,\\
 +\infty,&p\le1,
 \end{cases}
\qquad
 \lim_{\eta\downarrow0}\int_\eta^1 \frac{dx}{x^p}
 =
 \begin{cases}
 \dfrac{1}{1-p},&p<1,\\
 +\infty,&p\ge1.
 \end{cases}
\]
无穷远与零点的临界指数方向相反。
\end{example}

\section{Cauchy 准则}

\begin{theorem}[无穷区间的 Cauchy 准则]
\label{thm:improper-cauchy-infinity}
\(\int_a^\infty f\) 收敛，当且仅当对任意 \(\varepsilon>0\)，存在
\(A\ge a\)，使任意 \(v>u\ge A\) 都有
\[
 \abs{\int_u^v f(x)\,dx}<\varepsilon.
\]
\end{theorem}

\begin{proof}
令 \(F(R)=\int_a^R f\)。题中条件正是函数 \(F(R)\) 在
\(R\to\infty\) 时的 Cauchy 条件，因为
\[
 F(v)-F(u)=\int_u^v f.
\]
由函数极限的 Cauchy 准则与实数完备性即得。
\end{proof}

\begin{theorem}[有限奇点的 Cauchy 准则]
若 \(b\) 是右端奇点，则 \(\int_a^b f\) 收敛，当且仅当对任意
\(\varepsilon>0\)，存在 \(\delta>0\)，使
\[
 b-\delta<u<v<b
 \quad\Longrightarrow\quad
 \abs{\int_u^v f}<\varepsilon.
\]
\end{theorem}

\begin{proof}
令
\[
 G(\eta)=\int_a^{b-\eta}f(x)\,dx,\qquad \eta>0.
\]
当 \(\eta,\theta\downarrow0\) 时，\(G(\eta)-G(\theta)\) 正好是靠近
\(b\) 的一个短区间上的积分，差一个由端点次序决定的符号。因此题设条件
等价于 \(G(\eta)\) 在 \(\eta\downarrow0\) 时满足函数极限的 Cauchy
准则。实数完备性给出结论。
\end{proof}

\begin{proposition}[截断点的无关性]
\label{prop:improper-splitting-point}
若 \(f\) 在每个有限闭区间上可积，\(c,d\) 是两个有限截断点，则
\[
 \int_c^\infty f\text{ 收敛}
 \quad\Longleftrightarrow\quad
 \int_d^\infty f\text{ 收敛}.
\]
收敛时两者相差有限积分 \(\int_c^d f\)。双无穷区间定义中的分点也不影响
收敛性和积分值。
\end{proposition}

\begin{proof}
不妨设 \(c<d\)。对每个 \(R>d\)，区间可加性给出
\[
 \int_c^R f=\int_c^d f+\int_d^R f.
\]
第一项是与 \(R\) 无关的实数，所以两端截断积分同时有有限极限。双无穷
区间的情形分别在左右两侧应用同一论证。
\end{proof}

Cauchy 准则把收敛归结为尾积分的控制。上面的命题则说明，反常积分只在
奇点或无穷端附近增加了新的收敛问题；远离这些位置的有限一段不会改变
敛散性。

\section{绝对收敛、条件收敛与比较判别}

\begin{definition}
若 \(\int_a^\infty\abs f\) 收敛，则称
\(\int_a^\infty f\) \term{绝对收敛}。积分收敛但不绝对收敛时称为
\term{条件收敛}。
\end{definition}

\begin{theorem}[绝对收敛蕴含收敛]
\label{thm:improper-absolute-convergence}
若 \(\int_a^\infty\abs f<\infty\)，则
\(\int_a^\infty f\) 收敛，并且
\[
 \abs{\int_a^\infty f}
 \le\int_a^\infty\abs f.
\]
\end{theorem}

\begin{proof}
对任意 \(v>u\)，
\[
 \abs{\int_u^v f}\le\int_u^v\abs f.
\]
右端由 \(\int\abs f\) 的 Cauchy 准则在尾部任意小，因此
\(\int f\) 也满足 Cauchy 准则。对有限截断使用积分不等式，再令端点
趋于无穷即得最后一式。
\end{proof}

\begin{theorem}[比较判别]
\label{thm:improper-comparison}
设尾部有 \(0\le f\le g\)。若 \(\int g\) 收敛，则 \(\int f\)
收敛；若 \(\int f\) 发散到 \(+\infty\)，则 \(\int g\) 也发散到
\(+\infty\)。
\end{theorem}

\begin{proof}
取 \(A\) 使 \(0\le f(x)\le g(x)\) 对 \(x\ge A\) 成立。对 \(R\ge A\)，
\[
 0\le\int_A^R f\le\int_A^R g.
\]
正函数的两个截断积分都随 \(R\) 单调不减。若 \(\int_A^\infty g\)
收敛，则左端部分积分有上界，故由单调有界收敛而有有限极限。反之，若
\(\int_A^R f\to+\infty\)，同一不等式迫使 \(\int_A^R g\to+\infty\)。
最后用命题 \ref{prop:improper-splitting-point} 把截断点从 \(A\) 换回
原来的端点。
\end{proof}

\begin{corollary}[极限比较]
若 \(f,g>0\) 于尾部，且
\[
 \lim_{x\to\infty}\frac{f(x)}{g(x)}=L\in(0,\infty),
\]
则 \(\int f\) 与 \(\int g\) 同敛散。
\end{corollary}

\begin{proof}
尾部有 \(L/2\le f/g\le 3L/2\)，双向比较即可。
\end{proof}

若极限是零，只能由 \(\int g\) 收敛推出 \(\int f\) 收敛；若极限为
\(+\infty\)，只能得到相反方向。临界情形必须另作判断。

比较判别只利用绝对大小，因而特别适合非负函数和绝对值积分。对于带符号
函数，直接比较 \(f\le g\) 通常没有意义：两个函数的正负抵消可能完全
不同。此时或者先研究 \(\abs f\)，或者使用下一节的振荡判别。

\section{Dirichlet 与 Abel 判别}

正项比较不能解释振荡产生的抵消。下述证明把振荡部分的原函数和缓慢变化
部分分开。

\begin{theorem}[反常积分的 Dirichlet 判别]
\label{thm:improper-dirichlet}
设 \(f\) 在每个有限区间上连续，并且
\[
 F(x)=\int_a^x f(t)\,dt
\]
在 \([a,\infty)\) 上有界。设 \(g\in C^1([a,\infty))\) 单调且
\(g(x)\to0\)。则 \(\int_a^\infty f(x)g(x)\,dx\) 收敛。
\end{theorem}

\begin{proof}
不妨设 \(g\) 单调不增；另一情形对 \(-g\) 处理。由 \(g\to0\)，可在
有限起点后设 \(g\ge0\)。对 \(v>u\)，分部积分给出
\[
 \int_u^v fg
 =F(v)g(v)-F(u)g(u)-\int_u^v F(x)g'(x)\,dx.
\]
若 \(\abs F\le M\)，则因 \(g'\le0\)，
\[
 \abs{\int_u^v fg}
 \le M g(v)+M g(u)+M\int_u^v(-g')
 =2M g(u).
\]
当 \(u\to\infty\) 时右端趋于零，Cauchy 准则给出收敛。
\end{proof}

证明实际上给出了尾积分估计
\[
 \abs{\int_u^\infty f(x)g(x)\,dx}\le 2M\abs{g(u)}
\]
（在 \(g\) 最终单调并固定符号的尾部成立）。这条估计比“积分收敛”
提供了更多信息：若知道 \(g\) 的衰减速度，就同时得到余项的数量级。

\begin{example}
\[
 \int_1^\infty\frac{\sin x}{x}\,dx
\]
由 Dirichlet 判别收敛，因为 \(\int_1^R\sin x\,dx\) 有界而 \(1/x\)
单调趋零。它不绝对收敛；可在每个
\([k\pi+\pi/6,k\pi+5\pi/6]\) 上用
\(\abs{\sin x}\ge1/2\) 得到下界 \(1/(3(k+1))\)（\(k\ge1\)）。
对前 \(N\) 段相加，并用
\[
 \sum_{k=1}^N\frac1{k+1}\ge\int_2^{N+2}\frac{dt}{t}
 =\log\frac{N+2}{2}\longrightarrow+\infty,
\]
可见绝对值的截断积分无界。这里比较的是有限和与已经建立的积分，
不需要预先使用级数判别法。
\end{example}

\begin{theorem}[反常积分的 Abel 判别]
\label{thm:improper-abel}
设 \(f\) 在每个有限区间上连续。若 \(\int_a^\infty f\) 收敛，
\(g\in C^1([a,\infty))\) 单调且有界，
则 \(\int_a^\infty fg\) 收敛。
\end{theorem}

\begin{proof}
单调有界函数有有限极限 \(L\)。写
\[
 fg=L f+f(g-L).
\]
第一项的反常积分收敛；函数 \(g-L\) 单调趋零，而
\(F(x)=\int_a^x f\) 在无穷远有有限极限，又在每个有限闭区间连续，
因而在整个半轴有界。所以第二项
由 Dirichlet 判别收敛。
\end{proof}

\section{反常积分的变量代换与分部积分}

有限区间公式可以对每个截断积分使用，但必须再检查端点极限。

\begin{theorem}[反常积分中的变量代换]
\label{thm:improper-substitution}
设 \(\phi:[\alpha,\beta)\to[a,\infty)\) 连续，在
\((\alpha,\beta)\) 上属于 \(C^1\)，并且严格递增、满射，满足
\(\phi(\alpha)=a\) 和 \(\phi(t)\to\infty\)（\(t\to\beta-\)）。
若 \(f\) 在每个有限子区间连续，则
\[
 \int_a^\infty f(x)\,dx
 \quad\text{与}\quad
 \int_\alpha^\beta f(\phi(t))\phi'(t)\,dt
\]
同敛散，并在收敛时取相同值。若变换后的函数在 \(\alpha\) 附近无界，
新积分在该端也按反常积分理解；这一端的绝对收敛由上述假设自动保证。
\end{theorem}

\begin{proof}
固定 \(\alpha<r<s<\beta\)。在 \([r,s]\) 上可用有限区间代换公式，
得到
\[
 \int_r^s f(\phi(t))\phi'(t)\,dt
 =\int_{\phi(r)}^{\phi(s)}f(x)\,dx.
\]
\(\phi\) 递增，故 \(\phi'\ge0\)。对连续函数 \(\abs f\) 同样换元，
再令 \(r\downarrow\alpha\)，右端趋于有限积分
\(\int_a^{\phi(s)}\abs f\)。所以新下端绝对收敛，而且
\[
 \int_\alpha^s f(\phi(t))\phi'(t)\,dt
 =\int_a^{\phi(s)}f(x)\,dx.
\]
当 \(s\to\beta-\) 时，\(\phi(s)\to\infty\)。反过来，任意充分大的
\(R\) 都可唯一写为 \(\phi(s)\)，且 \(R\to\infty\) 必有
\(s\to\beta-\)：否则 \(s\) 会有一个小于 \(\beta\) 的上界，
从而 \(\phi(s)\) 也有界。故两侧上端极限同时存在且相等。
\end{proof}

变量代换并不能把一个发散积分自动变成收敛积分。定理中的“同敛散”来自
截断积分的逐项恒等，而不是来自形式上的微分记号
\(dx=\phi'(t)\,dt\)。若映射不是单调满射，或者新端点没有对应到原奇点，
就必须把区间分段并重新核对每一段的极限。

\begin{theorem}[反常分部积分]
\label{thm:improper-integration-by-parts}
设 \(u,v\in C^1([a,\infty))\)。若
\[
 \lim_{R\to\infty}u(R)v(R)
\]
存在且有限，并且 \(\int_a^\infty u'v\) 收敛，则
\(\int_a^\infty uv'\) 收敛，且
\[
 \int_a^\infty uv'
 =\lim_{R\to\infty}u(R)v(R)-u(a)v(a)
  -\int_a^\infty u'v.
\]
交换两项角色的版本同样成立。
\end{theorem}

\begin{proof}
在 \([a,R]\) 上使用分部积分公式，再令 \(R\to\infty\)。题设恰好保证
右端每一项都有有限极限。
\end{proof}

若边界项没有极限，不能只写形式上的“无穷端点代入”。同样，内部奇点的
变量代换必须分别作用于奇点两侧，不能用一个对称截断掩盖两段发散。

\section{Cauchy 主值}

通常反常积分要求每个无穷端点或奇点一侧的截断极限分别存在。某些带抵消
的积分还会采用对称截断；所得极限称为 Cauchy 主值，它是另一种极限。

\begin{definition}[无穷区间上的 Cauchy 主值]
设 \(f\) 在每个有限区间上可积。若极限
\[
 \lim_{R\to\infty}\int_{-R}^{R}f(x)\,dx
\]
存在且有限，则称它为 \(f\) 在全实线上的 \term{Cauchy 主值}，记作
\[
 \operatorname{p.v.}\!\int_{-\infty}^{\infty}f(x)\,dx.
\]
\end{definition}

\begin{definition}[内部奇点处的 Cauchy 主值]
设 \(a<c<b\)，且 \(f\) 在 \([a,c)\) 与 \((c,b]\) 的每个闭子区间
上可积。若
\[
 \lim_{\varepsilon\downarrow0}
 \left(\int_a^{c-\varepsilon}f(x)\,dx
       +\int_{c+\varepsilon}^b f(x)\,dx\right)
\]
存在且有限，则称该极限为 \(f\) 在 \(c\) 处的 Cauchy 主值。
\end{definition}

若通常反常积分收敛，则相应主值必存在且数值相同。这只需把两个分别收敛
的截断积分沿同一参数相加。反命题不成立，因为主值允许左右两侧的发散量
在取极限以前抵消。

\begin{example}[无穷远处的对称抵消]
函数 \(f(x)=x/(1+x^2)\) 为奇函数，所以
\[
 \int_{-R}^{R}\frac{x}{1+x^2}\,dx=0
\]
对每个 \(R>0\) 都成立，主值为零。然而
\[
 \int_0^R\frac{x}{1+x^2}\,dx=\frac12\log(1+R^2)\longrightarrow+\infty.
\]
左右两侧的通常反常积分均不收敛，因而全实线上的通常反常积分不存在。
\end{example}

\begin{example}[内部奇点与截断方式]
在 \([-1,1]\) 上，
\[
 \int_{-1}^{-\varepsilon}\frac{dx}{x}
 +\int_{\varepsilon}^{1}\frac{dx}{x}=0,
\]
所以 \(1/x\) 在零点处的主值为零；但两侧积分分别趋于
\(-\infty\) 与 \(+\infty\)。而采用不对称截断时，\(\alpha,\beta>0\)
给出
\[
 \int_{-1}^{-\alpha\varepsilon}\frac{dx}{x}
 +\int_{\beta\varepsilon}^{1}\frac{dx}{x}
 =\log\frac{\alpha}{\beta}.
\]
因此主值中的对称方式是定义的一部分，不能把不同截断方式得到的抵消结果
混同为通常积分。
\end{example}

\section*{习题}
\addcontentsline{toc}{section}{习题}

\subsection*{A组　定义与基本判别}
\begin{enumerate}[label=\textbf{\thechapter.\arabic*.}]
 \exerciseitem{improper-p-integrals-complete} 完整计算两类 \(p\)-型积分的截断
 极限，包括 \(p=1\) 的对数情形。
 \exerciseitem{improper-split-independence} 证明若
 \(\int_{-\infty}^\infty f\) 按定义收敛，则分割点 \(c\) 的选择不影响积分值。
 \exerciseitem{improper-principal-value-not-integral} 求
 \(f(x)=x/(1+x^2)\) 在全实线上的对称主值，并证明通常反常积分发散。
 \exerciseitem{improper-interior-singularity-split} 讨论
 \(\int_{-1}^1x^{-1}\,dx\) 的两侧积分和 Cauchy 主值，说明两种概念不同。
 \exerciseitem{improper-cauchy-tail-sequence} 证明无穷区间 Cauchy 准则等价于：
 对任意趋于 \(+\infty\) 的 \(u_n<v_n\)，都有
 \(\int_{u_n}^{v_n}f\to0\)。
\end{enumerate}

\subsection*{B组　比较与振荡}
\begin{enumerate}[label=\textbf{\thechapter.\arabic*.},resume]
 \exerciseitem{improper-log-power-infinity} 设 \(p,q\in\R\)，判定
 \[
  \int_e^\infty\frac{dx}{x^p(\log x)^q}
 \]
 的敛散性，完整处理 \(p=1\) 的临界情形。
 \exerciseitem{improper-log-power-zero} 设 \(p,q\in\R\)，判定
 \[
  \int_0^{1/2}\frac{dx}{x^p\abs{\log x}^{\,q}}
 \]
 的敛散性。
 \exerciseitem{improper-sine-over-power} 对 \(\alpha\in\R\)，判定
 \(\int_1^\infty \sin x/x^\alpha\,dx\) 的收敛、绝对收敛和发散范围。
 \exerciseitem{improper-dirichlet-quantitative-tail} 在 Dirichlet 判别的条件下
 证明尾积分估计
 \[
  \abs{\int_u^\infty fg}\le 2M\abs{g(u)}
 \]
 （在单调趋零的固定符号尾部）。
 \exerciseitem{improper-abel-from-dirichlet} 在本章 Abel 判别的假设下，
 不直接引用其结论，按
 \(g=L+(g-L)\) 的分解逐项验证全部条件。
 \exerciseitem{improper-absolute-not-necessary} 证明
 \(\int_1^\infty\sin x/x\) 条件收敛。
\end{enumerate}

\subsection*{C组　变换与反例}
\begin{enumerate}[label=\textbf{\thechapter.\arabic*.},resume]
 \exerciseitem{improper-substitution-inversion} 设 \(f:(0,1]\to\R\) 连续。
 用 \(x=1/t\) 把零点附近的
 反常积分化为无穷远积分，并重新得到 \(p\)-型判别。
 \exerciseitem{improper-gaussian-convergence} 不计算具体值，只用比较判别证明
 \(\int_0^\infty e^{-x^2}\,dx\) 收敛。
 \exerciseitem{improper-integration-by-parts-boundary} 构造 \(u,v\) 使
 \(\int u'v\) 收敛但 \(u(R)v(R)\) 无极限，从而说明反常分部积分的边界项
 条件不能删去。
 \exerciseitem{improper-conditional-substitution-warning} 证明
 \(\int_1^\infty\sin x/x\) 在满足定理 \ref{thm:improper-substitution}
 假设的代换下保持其值；说明为什么
 任意重排有限区间的贡献不属于变量代换定理。
 \optionalexerciseitem{improper-fresnel-dirichlet} 用代换 \(t=x^2\) 和
 Dirichlet 判别证明 \(\int_0^\infty\sin(x^2)\,dx\) 收敛但不绝对收敛。
\end{enumerate}
